\section{Fundamentos para resolver EDOs}

\subsection{¿Por qué se necesitan resolutores rápidos de EDOs}

El proceso de muestreo de los modelos de difusión es equivalente a resolver la PF-EDO desde $t=T$ (ruido puro) hasta $t=0$ (datos limpios). Los métodos ingenuos —como el muestreo de 1000 pasos en DDPM— son esencialmente una discretización de Euler-Maruyama para las SDEs, requiriendo una propagación hacia adelante de la red neuronal en cada paso. En aplicaciones prácticas, deseamos comprimir el número de pasos desde cientos a 10$\sim$20, manteniendo al mismo tiempo la calidad de la imagen.

\begin{importantbox}{Desafío central: muestreo de alta calidad con pocos pasos}
Cada "paso" del resolutor de EDOs requiere llamar una vez a la red de predicción de ruido $\boldsymbol{\epsilon}_\theta$, lo que constituye el cuello de botella computacional. Por tanto, reducir el número de pasos de resolución = reducir las inferencias de la red = acelerar la generación. El objetivo es generar imágenes de alta calidad dentro de 10$\sim$20 evaluaciones de función (NFE, Number of Function Evaluations).
\end{importantbox}

\subsection{Método de Euler}

Dada la EDO $\frac{d\mathbf{x}}{dt} = \mathbf{v}(\mathbf{x}, t)$, el método de Euler es la estrategia de integración numérica más sencilla:
$$\mathbf{x}_{t+h} = \mathbf{x}_t + h \cdot \mathbf{v}(\mathbf{x}_t, t)$$

Donde $h$ es el tamaño del paso. Intuitivamente, el método de Euler asume que el campo de velocidades $\mathbf{v}$ permanece constante en el intervalo $[t, t+h]$, usando la velocidad en el punto inicial para realizar una extrapolación lineal.

\begin{knowledgebox}{Intuición geométrica del método de Euler}
Imagina una curva en un plano bidimensional (la trayectoria real de la EDO). El método de Euler calcula la dirección tangente (campo de velocidades) en el punto actual y luego avanza un pequeño paso a lo largo de dicha tangente. A menor tamaño del paso, mayor precisión de la aproximación, pero se requieren más pasos. Cuando el tamaño del paso es grande, la aproximación lineal se desvía de la curva real, generando un error de truncamiento (truncation error). Para los modelos de difusión, la fórmula de muestreo de DDIM es en realidad una discretización de Euler para la PF-EDO.
\end{knowledgebox}

\subsection{Visión general de resolutores de EDO de orden superior}

Para mantener la precisión con tamaños de paso mayores, el campo del análisis numérico ha desarrollado una gran cantidad de resolutores de orden superior:

\begin{table}[H]
\centering
\begin{tabular}{lcc}
\toprule
\textbf{Resolutor} & \textbf{Orden} & \textbf{NFE por paso} \\
\midrule
Euler (método de Euler) & 1 & 1 \\
Heun (método de Euler mejorado/método de los trapecios) & 2 & 2 \\
Runge-Kutta 4 (RK4) & 4 & 4 \\
Dormand-Prince (DOPRI5) & 5 & 6 \\
\bottomrule
\end{tabular}
\caption{Resolutores de EDO genéricos comunes y sus características}
\end{table}

\begin{warningbox}{Limitaciones de los resolutores genéricos}
Los resolutores mencionados anteriormente son aplicables a \textbf{cualquier} EDO, pero no aprovechan la estructura especial de la PF-EDO. Por ejemplo, el Runge-Kutta 4 requiere 4 evaluaciones de red por paso, lo cual es costoso en escenarios de modelos de difusión. La innovación central de DPM-Solver radica en: aprovechar la estructura semilineal de la PF-EDO para lograr una precisión de orden superior bajo la premisa de realizar \textbf{solo 1 evaluación de red por paso}.
\end{warningbox}

\subsection{Resumen de la sección}

El método de Euler es la estrategia más básica para resolver EDOs y también es la forma equivalente al muestreo DDIM. Los resolutores de orden superior genéricos (como RK4) pueden mejorar la precisión, pero requieren múltiples evaluaciones de red por paso. El enfoque principal de esta conferencia posterior será cómo aprovechar la estructura especial de las EDOs de los modelos de difusión para diseñar resolutores de orden superior del tipo "almuerzo gratis".
